\chapter*{}
\noindent{\textbf{RESUMO}}

\noindent{    
    Neste trabalho é feito um teste de hipótese, uma análise para descobrir quais os 
    impactos na escolha da arquitetura, paradigmas e a linguagem para desenvolver um jogo. 
    Utilizando \textit{Tetris} como estudo de caso e selecionando 4 tecnologias amplamente utilizadas (\textit{Godot}, Java com a biblioteca LibGDX, C++ com a biblioteca Raylib e Python com Pygame),
    observando em cada uma das tecnologias a facilidade de aprendizado, modularidade, estrutura de desenvolvimento, vantagens e limitações.
    A metodologia se baseia em uma abordagem prática, aplicando recursos de cada tecnologia de acordo com que cada jogo é implementado, utilizando documentações oficiais e materias complementares,tendo como resultado a implementação dos quatro jogos \textit{Tetris}, observou-se que as tecnologias mais acessíveis e práticas foram a \textit{Godot} e o Python com Pygame, enquanto Java com LibGDX e C++ com Raylib apresentaram maior complexidade de implementação.
    }

\vspace{\onelineskip}

% todas em letras minúsculas, separadas por ponto e vírgula (;)
\noindent{\textbf{Palavras-chave}: desenvolvimento de jogos; paradigmas; comparativo; tetris; arquitetura; java; C++; python; godot;}

